
\textbf{Lemme 4.2.} \uline{Sous les conditions de} 4.1, @@M0@@ \uline{est un morphisme de descente pour la catégorie des faisceaux localement libres~; c'est-à-dire pour tout couple} @@M1@@ \uline{de faisceaux localement libres sur} @@M2@@\uline{, le diagramme}
@@M3@@
\uline{est exact.}

En effet, si @@M4@@, on doit montrer que la suite
@@M5@@
est exacte, ou bien (d'après EGA @@M6@@ 5.4.10.1) que la suite
@@M7@@
est exacte. Mais la suite (*) @@M8@@ est exacte par hypothèse~; d'où le lemme, en lui appliquant le foncteur @@M9@@, qui est exact à gauche, @@M10@@ étant plat sur @@M11@@.

Soit @@M12@@ un faisceau inversible sur @@M13@@ muni d'une rigidification @@M14@@ relativement à @@M15@@~; d'après 4.2 on a une équivalence entre la catégorie @@M16@@ des couples @@M17@@ d'un faisceau inversible @@M18@@ sur @@M19@@ et d'un isomorphisme @@M20@@, et la catégorie @@M21@@ des données de descente effectives sur @@M22@@ relativement à @@M23@@. Mais si @@M24@@, comme @@M25@@ est rigidifié relativement à @@M26@@, @@M27@@ est muni canoniquement d'une rigidification relativement à @@M28@@. Par conséquent @@M29@@ est équivalente à la catégorie discrète définie par l'ensemble des classes de faisceaux inversibles @@M30@@-rigidifiés @@M31@@ sur @@M32@@, tels que @@M33@@. Mais c'est précisément l'ensemble @@M34@@ de @@M35@@-points de l'image réciproque dans @@M36@@ du @@M37@@-point de @@M38@@ défini par @@M39@@. De plus cette description de @@M40@@ est fonctorielle en @@M41@@. Donc @@M42@@ s'identifie au foncteur @@M43@@, et on peut oublier les rigidifications par la suite.

